% =========================================================
% Irrationals as a Subset of the Real Line
% =========================================================

\subsection{Irrationals as a Subset of the Real Line}
\label{sec:irrationals-as-a-subset}

\begin{exposition}
The irrational numbers form a large and unavoidable subset of the real line.
They lie in every nonempty interval, they approximate and are approximated by
rational numbers, and they inherit some but not all of the structural features
of $\mathbb{R}$.
\end{exposition}

\begin{theorembox}{Theorem (Density of $\mathbb{Q}$ in $\mathbb{R}$)}
\begin{theorem}[Density of $\mathbb{Q}$ in $\mathbb{R}$]
\label{thm:q-dense-in-r}
If $x,y\in\mathbb{R}$ and $x<y$, then there exists $q\in\mathbb{Q}$ such that
\[
x<q<y.
\]
\end{theorem}
\end{theorembox}
\begin{dependencies}
\hyperref[thm:archimedean-property]{Archimedean property of $\mathbb{R}$},
\hyperref[def:rationals]{definition of $\mathbb{Q}$ as ratios of integers}.
Missing formal dependency target: the floor/integer-part lemma for real numbers.
\end{dependencies}
\begin{remark*}[Proof idea]
Scale the interval $(x,y)$ until its length exceeds $1$, place an integer
between the scaled endpoints, and divide back down.
\end{remark*}
\begin{remark*}[Proof]
\hyperref[prf:q-dense-in-r]{Go to proof.}
\end{remark*}

\begin{corollarybox}{Corollary (Infinitely Many Rational Numbers Lie Between Two Distinct Reals)}
\begin{corollary}[Infinitely Many Rational Numbers Lie Between Two Distinct Reals]
\label{cor:infinitely-many-rationals-between}
Between any two distinct real numbers there are infinitely many rational numbers.
\end{corollary}
\end{corollarybox}
\begin{dependencies}
\hyperref[thm:q-dense-in-r]{Density of $\mathbb{Q}$ in $\mathbb{R}$}.
\end{dependencies}
\begin{remark*}[Proof idea]
After finding one rational in the interval, apply density again inside the
smaller subintervals to continue producing new ones.
\end{remark*}
\begin{remark*}[Proof]
\hyperref[prf:infinitely-many-rationals-between]{Go to proof.}
\end{remark*}

\begin{theorembox}{Theorem (Between Any Two Real Numbers Lies an Irrational)}
\begin{theorem}[Between Any Two Real Numbers Lies an Irrational]
\label{thm:between-any-two-reals-is-an-irrational}
If $x,y\in\mathbb{R}$ and $x<y$, then there exists
$\alpha\in\mathbb{R}\setminus\mathbb{Q}$ such that
\[
x<\alpha<y.
\]
\end{theorem}
\end{theorembox}
\begin{dependencies}
\hyperref[thm:q-dense-in-r]{Density of $\mathbb{Q}$ in $\mathbb{R}$},
\hyperref[thm:rational-plus-irrational-is-irrational]{Theorem~\ref*{thm:rational-plus-irrational-is-irrational}},
and the irrationality of $\sqrt{2}$.
\end{dependencies}
\begin{remark*}[Proof idea]
Choose a sufficiently small rational perturbation $q$ and shift it by a small
irrational quantity such as $\sqrt{2}/n$.
\end{remark*}
\begin{remark*}[Proof]
\hyperref[prf:between-any-two-reals-is-an-irrational]{Go to proof.}
\end{remark*}

\begin{corollarybox}{Corollary (The Irrationals Are Dense in $\mathbb{R}$)}
\begin{corollary}[The Irrationals Are Dense in $\mathbb{R}$]
\label{cor:irrationals-dense-in-r}
Every nonempty open interval in $\mathbb{R}$ contains an irrational number.
\end{corollary}
\end{corollarybox}
\begin{dependencies}
\hyperref[thm:between-any-two-reals-is-an-irrational]{Theorem~\ref*{thm:between-any-two-reals-is-an-irrational}}.
\end{dependencies}
\begin{remark*}[Proof idea]
Unpack what the theorem says for an arbitrary open interval $(a,b)$.
\end{remark*}
\begin{remark*}[Proof]
\hyperref[prf:irrationals-dense-in-r]{Go to proof.}
\end{remark*}

\begin{corollarybox}{Corollary (Every Irrational Can Be Approximated by Rationals)}
\begin{corollary}[Every Irrational Can Be Approximated by Rationals]
\label{cor:every-irrational-can-be-approximated-by-rationals}
For every $\alpha\in\mathbb{R}\setminus\mathbb{Q}$ and every $\varepsilon>0$,
there exists $q\in\mathbb{Q}$ such that
\[
|\alpha-q|<\varepsilon.
\]
\end{corollary}
\end{corollarybox}
\begin{dependencies}
\hyperref[thm:q-dense-in-r]{Density of $\mathbb{Q}$ in $\mathbb{R}$}.
\end{dependencies}
\begin{remark*}[Proof idea]
Apply rational density to $(\alpha-\varepsilon,\alpha+\varepsilon)$.
\end{remark*}
\begin{remark*}[Proof]
\hyperref[prf:every-irrational-can-be-approximated-by-rationals]{Go to proof.}
\end{remark*}

\begin{corollarybox}{Corollary (Every Rational Can Be Approximated by Irrationals)}
\begin{corollary}[Every Rational Can Be Approximated by Irrationals]
\label{cor:every-rational-can-be-approximated-by-irrationals}
For every $q\in\mathbb{Q}$ and every $\varepsilon>0$, there exists
$\alpha\in\mathbb{R}\setminus\mathbb{Q}$ such that
\[
|\alpha-q|<\varepsilon.
\]
\end{corollary}
\end{corollarybox}
\begin{dependencies}
\hyperref[cor:irrationals-dense-in-r]{Corollary~\ref*{cor:irrationals-dense-in-r}}.
\end{dependencies}
\begin{remark*}[Proof idea]
Apply irrational density to $(q-\varepsilon,q+\varepsilon)$.
\end{remark*}
\begin{remark*}[Proof]
\hyperref[prf:every-rational-can-be-approximated-by-irrationals]{Go to proof.}
\end{remark*}

\begin{theorembox}{Theorem (Closure of $\mathbb{Q}$ in $\mathbb{R}$)}
\begin{theorem}[Closure of $\mathbb{Q}$ in $\mathbb{R}$]
\label{thm:closure-of-q-is-r}
The closure of $\mathbb{Q}$ in $\mathbb{R}$ is all of $\mathbb{R}$.
\end{theorem}
\end{theorembox}
\begin{dependencies}
\hyperref[thm:q-dense-in-r]{Density of $\mathbb{Q}$ in $\mathbb{R}$}.
\end{dependencies}
\begin{remark*}[Proof idea]
Every real point can be approximated arbitrarily closely by rationals, so every
real point lies in the closure of $\mathbb{Q}$.
\end{remark*}
\begin{remark*}[Proof]
\hyperref[prf:closure-of-q-is-r]{Go to proof.}
\end{remark*}

\begin{corollarybox}{Corollary ($\mathbb{Q}$ Is Not Closed in $\mathbb{R}$)}
\begin{corollary}[$\mathbb{Q}$ Is Not Closed in $\mathbb{R}$]
\label{cor:q-not-closed-in-r}
The set $\mathbb{Q}$ is not closed as a subset of $\mathbb{R}$.
\end{corollary}
\end{corollarybox}
\begin{dependencies}
\hyperref[thm:closure-of-q-is-r]{Theorem~\ref*{thm:closure-of-q-is-r}}.
\end{dependencies}
\begin{remark*}[Proof idea]
Its closure is all of $\mathbb{R}$, which is strictly larger than $\mathbb{Q}$.
\end{remark*}
\begin{remark*}[Proof]
\hyperref[prf:q-not-closed-in-r]{Go to proof.}
\end{remark*}

\begin{corollarybox}{Corollary ($\mathbb{Q}$ Is Not Open in $\mathbb{R}$)}
\begin{corollary}[$\mathbb{Q}$ Is Not Open in $\mathbb{R}$]
\label{cor:q-not-open-in-r}
The set $\mathbb{Q}$ is not open as a subset of $\mathbb{R}$.
\end{corollary}
\end{corollarybox}
\begin{dependencies}
\hyperref[cor:irrationals-dense-in-r]{Corollary~\ref*{cor:irrationals-dense-in-r}}.
\end{dependencies}
\begin{remark*}[Proof idea]
Every open interval around a rational point contains irrational points.
\end{remark*}
\begin{remark*}[Proof]
\hyperref[prf:q-not-open-in-r]{Go to proof.}
\end{remark*}

\begin{corollarybox}{Corollary ($\mathbb{Q}$ Is Not Discrete in $\mathbb{R}$)}
\begin{corollary}[$\mathbb{Q}$ Is Not Discrete in $\mathbb{R}$]
\label{cor:q-not-discrete-in-r}
The set $\mathbb{Q}$ is not discrete as a subset of $\mathbb{R}$.
\end{corollary}
\end{corollarybox}
\begin{dependencies}
\hyperref[thm:between-any-two-rationals-is-a-rational]{Theorem~\ref*{thm:between-any-two-rationals-is-a-rational}}.
\end{dependencies}
\begin{remark*}[Proof idea]
No rational point is isolated, because every interval around it contains other
rationals.
\end{remark*}
\begin{remark*}[Proof]
\hyperref[prf:q-not-discrete-in-r]{Go to proof.}
\end{remark*}

\begin{theorembox}{Theorem (The Irrationals Are Uncountable)}
\begin{theorem}[The Irrationals Are Uncountable]
\label{thm:irrationals-uncountable}
The set $\mathbb{R}\setminus\mathbb{Q}$ is uncountable.
\end{theorem}
\end{theorembox}
\begin{dependencies}
\hyperref[thm:R-uncountable]{$\mathbb{R}$ is uncountable},
\hyperref[thm:Q-countable]{$\mathbb{Q}$ is countable}.
Missing formal dependency target: removing a countable subset from an uncountable
set leaves an uncountable set.
\end{dependencies}
\begin{remark*}[Proof idea]
Subtracting a countable subset from an uncountable set leaves an uncountable
set.
\end{remark*}
\begin{remark*}[Proof]
\hyperref[prf:irrationals-uncountable]{Go to proof.}
\end{remark*}

\begin{theorembox}{Theorem (The Irrationals Are Not Complete as an Ordered Set)}
\begin{theorem}[The Irrationals Are Not Complete as an Ordered Set]
\label{thm:irrationals-not-complete-as-ordered-set}
The ordered set $\mathbb{R}\setminus\mathbb{Q}$ is not complete.
\end{theorem}
\end{theorembox}
\begin{dependencies}
\hyperref[def:least-upper-bound-property]{least upper bound property},
\hyperref[def:supremum]{definition of supremum},
\hyperref[def:rationals]{definition of $\mathbb{Q}$}.
Missing formal dependency target: completeness failure for an inherited ordered
subset whose real supremum lies outside the subset.
\end{dependencies}
\begin{remark*}[Proof idea]
Consider the set
\[
A=\{\alpha\in\mathbb{R}\setminus\mathbb{Q}:\alpha<1\}.
\]
Its least upper bound in $\mathbb{R}$ is $1$, but $1$ is not irrational.
\end{remark*}
\begin{remark*}[Proof]
\hyperref[prf:irrationals-not-complete-as-ordered-set]{Go to proof.}
\end{remark*}
